% G19 insertion used by the Rev5 manuscript source.

\begin{proposition}[Exact strong/weak-coupling separation]
\label{prop:g19-axis}
For $SU(3)$ the Hamiltonian coordinate of this paper is
$u=g_H^{-4}$.  Put
\[
 u_*:=\min\left\{\frac{c_1}{27},\frac{1}{891c_2}\right\},
\]
the coupling radius in Theorem~\ref{thm:riesz-island}.  If $a\downarrow0$ is
any asymptotically-free Hamiltonian trajectory with $g_H(a)\to0$, then
$u(a)\to\infty$.  In particular, there is $a_0>0$ such that
$u(a)>u_*$ for $0<a<a_0$.  No tail of a continuum trajectory therefore lies
in the proved Riesz-island domain.
\end{proposition}

\begin{proof}
The function $x\mapsto x^{-4}$ tends to infinity as $x\downarrow0$.
The eventual inequality is the definition of this limit.
\end{proof}

This is stronger than a warning about extrapolation: the hypotheses of the
fixed-cutoff theorem and of the continuum trajectory are eventually disjoint.
Numerical evaluation of the existential constants $c_1,c_2$ would not change
that fact.

\begin{conditional}[Necessary continuum scaling of an electric-unit branch]
\label{cond:g19-scaling}
Use the universal pure-gauge beta-function coefficients and the standard
two-loop lattice scaling law~\cite{GrossWilczek1973,Politzer1973,LuscherWeisz1995}.
Let
\[
 b_0=\frac{11N}{48\pi^2},\qquad
 b_1=\frac{34N^2}{3(16\pi^2)^2},\qquad
 p=\frac{b_1}{2b_0^2}=\frac{51}{121}.
\]
Assume that the Hamiltonian coupling has been matched to an
asymptotically-free Wilson scheme, including the anisotropy/time
normalisation, so that for a finite $\Lambda_H>0$
\begin{equation}\label{eq:g19-asymptotic-scaling}
 a\Lambda_H=
 \exp\!\left[-\frac{1}{2b_0g_H(a)^2}\right]
 \bigl(b_0g_H(a)^2\bigr)^{-p}[1+o(1)].
\end{equation}
Assume also that a vacuum-subtracted electric-unit excitation $E(u)$ obeys
\begin{equation}\label{eq:g19-hamiltonian-normalisation}
 aM(a)=c_Hg_H(a)^2E(u(a))[1+o(1)],\qquad c_H>0,
\end{equation}
and $M(a)\to M\in(0,\infty)$.  With $u=g_H^{-4}$, necessarily
\begin{equation}\label{eq:g19-necessary-law}
 E(u(a))=
 \frac{M}{c_H\Lambda_H}b_0^{-p}
 u(a)^{(1+p)/2}
 \exp\!\left[-\frac{\sqrt{u(a)}}{2b_0}\right][1+o(1)].
\end{equation}
For $SU(3)$ the essential exponential is
$\exp[-(8\pi^2/11)\sqrt u]$ and $(1+p)/2=86/121$.
\end{conditional}

The Hamer crosswalk printed earlier gives
$aM_{\rm phys}=(g^2/2)M_A(x)$ and
$E(u)=M_A(2u)/2$.  Hence $x=2u$ and $g=g_H$ would fix $c_H=1$,
not $1/2$: the half is already present in $E=M_A/2$.  We leave $c_H$
explicit because the upstream hostile audit still records the operator
normalisation $H=W/2$, $u=x/2$ as an asserted bridge and because the
Hamiltonian--Wilson anisotropy/time map needed at weak coupling has not been
derived.  Proving that map would set $c_H=1$ in the present convention.

\begin{proof}
Set $x=g_H^{-2}=\sqrt u$.  Equation
\eqref{eq:g19-asymptotic-scaling} gives
$a=\Lambda_H^{-1}e^{-x/(2b_0)}(x/b_0)^p[1+o(1)]$, while
\eqref{eq:g19-hamiltonian-normalisation} gives
$E=aMx/c_H[1+o(1)]$.  Substitution proves
\eqref{eq:g19-necessary-law}.
\end{proof}

\begin{corollary}[Finite Taylor data cannot supply the continuum mass]
\label{cor:g19-polynomial-obstruction}
Under the matching assumptions of Conditional result
\ref{cond:g19-scaling}, substituting any nonzero polynomial $P(u)$ for
$E(u)$ makes
\[
 \left|\frac{c_Hg_H(a)^2P(u(a))}{a}\right|\longrightarrow\infty.
\]
Thus no nonzero finite strong-coupling Taylor truncation yields a finite
continuum mass.
\end{corollary}

\begin{proof}
If $P$ has degree $d$ and leading coefficient $q_d\ne0$, then with
$x=\sqrt u$ the expression is
\[
 c_H\Lambda_Hq_db_0^p
 e^{x/(2b_0)}x^{2d-1-p}[1+o(1)].
\]
Its absolute value diverges because the exponential dominates every real
power.  The conclusion is about finite truncations, not about a
nonperturbatively continued branch.
\end{proof}

\begin{conditional}[Bare absolute-weight power counting]
\label{cond:g19-bare-residue}
Work first in a periodic box of fixed physical volume and let
\[
 \widetilde O_{a,n}=L(a)^{-3/2}\sum_x
       \operatorname{ImTr}U_{p_n(x)}
\]
be the volume-normalized zero-momentum source made from plaquettes normal to
direction $n$.  Suppose a continuum operator matching exists in the form
\begin{equation}
 \operatorname{ImTr}U_{p_n(x)}
 =a^d Z_d(a)\mathcal O_n(x)+o(a^dZ_d(a))
 \label{eq:g19-local-source-matching}
\end{equation}
in matrix elements against a fixed normalized state, and suppose the
corresponding continuum matrix element is finite.  Then
\begin{equation}
 \left|\langle\Psi,\widetilde O_{a,n}\Omega\rangle\right|^2
 =O\!\left(a^{2d-3}|Z_d(a)|^2\right).
 \label{eq:g19-overlap-power}
\end{equation}
For the charge-odd plaquette source, tracelessness and
\begin{equation}
 \operatorname{ImTr}e^{iX}
 =-\frac16\operatorname{Tr}X^3+O(\|X\|^5)
 \label{eq:g19-imtr-expansion}
\end{equation}
give $d=6$ and
$\mathcal O_n\propto d^{abc}B_n^aB_n^bB_n^c$.  Thus this state's absolute
spectral-atom weight for the unrenormalized, volume-normalized source has the
power $a^9|Z_6(a)|^2$.  If $Z_6$ contains only the usual logarithmic matching
factors, that absolute weight vanishes as $a\downarrow0$.
\end{conditional}

\begin{proof}
At fixed physical volume, $L(a)^3\asymp a^{-3}$.  Replacing the lattice sum
by its Riemann sum in~\eqref{eq:g19-local-source-matching} gives
\[
 L(a)^{-3/2}\sum_x a^dZ_d(a)\mathcal O_n(x)
 =a^{d-3/2}Z_d(a),
   \frac1{\sqrt{V}}\int_V\mathcal O_n(y)\,d^3y+o(a^{d-3/2}Z_d(a)).
\]
Squaring proves~\eqref{eq:g19-overlap-power}.  For a plaquette,
$X=a^2g_HF_{ij}+O(a^3)$; the linear term in
$\operatorname{Tr}e^{iX}$ vanishes and the first imaginary term is cubic,
which proves~\eqref{eq:g19-imtr-expansion} and $d=6$.
\end{proof}

This conditional result is an absolute-weight power count for the
volume-normalized operator $\widetilde O_{a,n}$ above; it is not a
statement about the Hilbert-normalized spectral fraction.  At fixed physical
Euclidean time the same overall operator normalization multiplies the bottom
atom and the full correlator, so it cancels from their ratio.  Consequently
the factor $a^9|Z_6(a)|^2$ alone proves neither collapse nor survival of a
normalized residue.  Smearing over a radius fixed in physical units changes
the operator and its matching factors; a uniform lower bound for the
renormalized source fraction still requires independent spectral-measure and
matching control.  Moreover, the three diagonal tensors
$d^{abc}B_n^aB_n^bB_n^c$ transform as $T_1$ under the cubic group but belong
to a symmetric rank-three continuum tensor; the label $T_1^{+-}$ alone does
not distinguish continuum spin one from spin three.  Auxiliary
$A_2^{+-}$ and $T_2^{+-}$ operators are required for that degeneracy test.
This is the standard cubic-subduction partner method used in continuum-spin
assignments~\cite{MorningstarPeardon1999,AthenodorouTeper2020}.

\begin{proposition}[Passage of a source-visible spectral atom]
\label{prop:g19-atom-passage}
Let $\mu_n$ be positive locally finite Borel measures on $[0,\infty)$,
and suppose
\[
 M_n=\inf\operatorname{supp}\mu_n\in[m_-,m_+],
 \qquad \mu_n(\{M_n\})\ge z_*>0 .
\]
Assume that the source spectral measures converge vaguely,
$\mu_n\Rightarrow\mu$.  After passage to any subsequence on which
$M_n\to M$, one has
\begin{equation}
 \operatorname{supp}\mu\subset[M,\infty),
 \qquad \mu(\{M\})\ge z_* .
 \label{eq:g19-atom-passage}
\end{equation}
Thus a uniformly source-visible bottom atom survives as a nonzero atom of
the limiting spectral measure; no isolation annulus, spin assignment, or
Haag--Ruelle hypothesis is needed for this weaker conclusion.
\end{proposition}

\begin{proof}
Compactness of $[m_-,m_+]$ supplies the stated subsequence.  If
$K\Subset[0,M)$, then $K\cap\operatorname{supp}\mu_n=\varnothing$ for all
sufficiently large $n$, and vague convergence gives $\mu(K)=0$.  Hence the
limiting support is contained in $[M,\infty)$.

For $\varepsilon>0$, choose $f_\varepsilon\in C_c([0,\infty))$ with
$0\le f_\varepsilon\le1$, equal to one on
$[M-\varepsilon/2,M+\varepsilon/2]\cap[0,\infty)$ and supported in
$[M-\varepsilon,M+\varepsilon]\cap[0,\infty)$.  Eventually
$f_\varepsilon(M_n)=1$, so positivity and vague convergence give
\[
 \int f_\varepsilon\,d\mu
 =\lim_n\int f_\varepsilon\,d\mu_n\ge z_* .
\]
Consequently
$\mu([M-\varepsilon,M+\varepsilon]\cap[0,\infty))\ge z_*$.  Continuity
from above as $\varepsilon\downarrow0$ proves
$\mu(\{M\})\ge z_*$.
\end{proof}

For a positive transfer Hamiltonian and a charge-odd source, write
\[
 C_n(t)=\int e^{-tE}\,d\mu_n(E),\qquad
 Z_n=\mu_n(\{M_n\}).
\]
Osterwalder--Schrader convergence~\cite{OsterwalderSeiler1978} strong enough
to imply vague convergence
of these source measures therefore turns a uniform lower bound on $Z_n$
and a two-sided physical bound on $M_n$ into the nonzero limiting atom in
Proposition~\ref{prop:g19-atom-passage}.  This is an exact reduction, not a
proof of either uniform bound or of the required OS convergence.

\begin{lemma}[What a finite-time correlator can certify]
\label{lem:g19-residue-bound}
Let a finite positive measure have the form
\[
 \mu=Z\delta_M+\nu,
 \qquad \operatorname{supp}\nu\subset[M+\Delta,\infty),
 \qquad \Delta>0,
\]
and put $C(t)=\int e^{-tE}\,d\mu(E)$ and $F(t)=e^{Mt}C(t)$.  Then
\begin{equation}
 Z\le F(t),
 \qquad
 Z\ge
 \frac{F(t)-e^{-\Delta t}C(0)}{1-e^{-\Delta t}} .
 \label{eq:g19-residue-bound}
\end{equation}
The second inequality may be replaced by zero when its right-hand side is
negative.
\end{lemma}

\begin{proof}
Positivity gives $F(t)\ge Z$.  The spectral gap gives
\[
 F(t)
 =Z+\int e^{-t(E-M)}\,d\nu(E)
 \le Z+e^{-\Delta t}\bigl(C(0)-Z\bigr),
\]
and rearrangement proves the lower bound.
\end{proof}

Without a gap above the atom, every finite-$t$ value of $F(t)$ is only an
upper bound on $Z$; an absolutely continuous threshold can imitate a
mass edge at finite time while having $Z=0$.  Lemma
\ref{lem:g19-residue-bound} therefore explains exactly where a uniform
external gap can enter a nonperturbative residue proof, while
Proposition~\ref{prop:g19-atom-passage} shows that isolation is not part of
the final atom-passage statement itself.

\begin{proposition}[Conditional OS gap passage]
\label{prop:g19-mosco}
Let $q_n$ be nonnegative closed forms of physical transfer Hamiltonians
$H_n$, with vacuum projections $P_n$, after isometric embedding in a common
Hilbert space.  Suppose $q_n$ Mosco-converges to the form $q$ of an
OS-reconstructed Hamiltonian $H$, $P_n\to P$ strongly,
$\operatorname{ran}P=\ker H$, and for one $m>0$ independent of $n$,
\[
 q_n[\psi]\ge m\lVert(1-P_n)\psi\rVert^2
 \quad(\psi\in D(q_n)).
\]
Then
\[
 q[\psi]\ge m\lVert(1-P)\psi\rVert^2,
 \qquad
 \operatorname{spec}(H)\subset\{0\}\cup[m,\infty).
\]
\end{proposition}

\begin{proof}
For $\psi\in D(q)$ choose a Mosco recovery sequence $\psi_n\to\psi$.
The recovery and liminf properties give $q_n[\psi_n]\to q[\psi]$.
Strong convergence of the projections gives
$(1-P_n)\psi_n\to(1-P)\psi$.  Passing to the limit proves the form bound;
the spectral inclusion follows from the spectral theorem.
\end{proof}

\begin{proposition}[Deterministic reflection-positive pushforward]
\label{prop:g19-rp-pushforward}
Let $(X,\mu,\theta;\mathcal A_+)$ be a reflected probability space with
\[
 (\Theta F)(U)=\overline{F(\theta U)},\qquad
 \int_X(\Theta F)F\,d\mu\ge0\quad(F\in\mathcal A_+).
\]
Let $(X',\theta';\mathcal A'_+)$ be a coarse reflected configuration space,
and let $P:X\to X'$ be measurable.  If
\begin{equation}
 P\theta=\theta'P\quad\mu\text{-a.e.},
 \qquad P^*\mathcal A'_+\subset\mathcal A_+,
 \label{eq:g19-rp-map}
\end{equation}
then the pushforward $P_\#\mu$ is reflection positive on
$\mathcal A'_+$.  On the gauge-invariant algebra the first equality may hold
only modulo a coarse gauge transformation.
\end{proposition}

\begin{proof}
For $F'\in\mathcal A'_+$ put $F=F'\circ P$.  Then
\[
 \int_{X'}(\Theta'F')F'\,d(P_\#\mu)
 =\int_X\overline{F'(\theta'P(U))}F'(P(U))\,d\mu(U)
 =\int_X(\Theta F)F\,d\mu\ge0.
\]
If equivariance holds modulo a coarse gauge transformation, gauge invariance
of $F'$ gives the same middle equality.
\end{proof}

\begin{corollary}[Reflection-adapted Balaban-form blocking]
\label{cor:g19-balaban-rp}
Let the Wilson measure live on a periodic lattice of even temporal extent and
be reflection positive on the gauge-invariant positive-time algebra.  Let an
odd block factor $L$ divide that extent.  Use centered $L^d$ blocks, put the
reflection plane between two adjacent layers of block centres, and choose the
reference trees and Balaban paths in reflection pairs, with reversed coarse
bonds defined by group inverse.  Then the raw deterministic Balaban-form
average of Eq.~(15) in Ref.~\cite{Balaban1985Averaging} pushes the Wilson
measure to a reflection-positive coarse measure.
\end{corollary}

\begin{proof}
Choose paths so that
\[
 \theta\Gamma_{c,x}=\Gamma_{\theta'c,\theta x},
 \qquad \Gamma_{c^{-1},x'}=\Gamma_{c,x}^{-1}.
\]
Holonomy is functorial under reflection, while logarithm and exponential
commute with unitary conjugation.  Reindexing the equal-weight average by
$x\mapsto\theta x$ proves $P(\theta U)=\theta'P(U)$ off the principal-log
branch set, which is Haar-null; Balaban's Eq.~(11) supplies gauge covariance.
For a positive coarse bond, both endpoint blocks lie wholly in the fine
positive half.  The locality statement following Eq.~(43) of
Ref.~\cite{Balaban1985Averaging} therefore gives
$P^*\mathcal A'_+\subset\mathcal A_+$.  Proposition
\ref{prop:g19-rp-pushforward} applies, using Wilson reflection positivity from
Ref.~\cite{OsterwalderSeiler1978}.
\end{proof}

The corollary proves preservation for the raw deterministic average with
reflection-adapted block and path data.  It does not cover Balaban's literal
lower-corner convention, field-dependent gauge fixing, an asymmetric
large/small-field partition, a softened stochastic constraint, or a complete
multistep effective action.  For the literal convention the finite-volume
geometry gate is now exactly
\begin{equation}
 P_{\rm CMP98}(\theta U)
 =g_\theta(U)\cdot\theta'P_{\rm CMP98}(U),
 \qquad
 \operatorname{supp}P_c\subset\Lambda_+
 \quad(c\in E'_+),
 \label{eq:g19-balaban-gate}
\end{equation}
where $g_\theta$ is the coarse gauge change induced by moving the reflected
upper-corner basepoint back to the chosen lower corner.  Block alignment by
itself does not prove Eq.~\eqref{eq:g19-balaban-gate}.

\begin{remark}[Equivariant Markov blocking is insufficient]
\label{rem:g19-markov-counterexample}
Let the fine space be one point.  On the coarse space
$X'=\{(s_-,s_+):s_\pm=\pm1\}$ with reflection
$(s_-,s_+)\mapsto(s_+,s_-)$, take the reflection-invariant law
\[
 \nu=\tfrac12\delta_{(+1,-1)}+\tfrac12\delta_{(-1,+1)}.
\]
The resulting Markov kernel is reflection equivariant and maps every
positive-time observable to a fine constant.  Nevertheless, for
$F(s_-,s_+)=s_+$,
\[
 \int(\Theta'F)F\,d\nu=\mathbb E_\nu[s_-s_+]=-1.
\]
Thus equivariance and one-sided mapping do not suffice for a stochastic
kernel.  A sufficient replacement is a reflection factorisation: there is a
one-sided $T:\mathcal A'_+\to\mathcal A_+$ such that
\[
 K\bigl((\Theta'F)G\bigr)=(\Theta TF)(TG)
 \qquad(F,G\in\mathcal A'_+).
\]
\end{remark}

The hypotheses of Proposition~\ref{prop:g19-mosco} are not supplied here.
Theorem~\ref{thm:riesz-island} controls a finite-volume cluster in a
fixed-cutoff small-$u$ domain, not the whole physical transfer Hamiltonian
along $a\downarrow0$.  The separate local weak-well expansion in the upstream
corpus also has no proved operator identity relating its $\beta_{\rm loc}$ to
$u$.  The exact next G19 object is therefore a Hamiltonian--Wilson matching
package: the physical energy prefactor, anisotropy/speed-of-light
renormalisation, $g_H$--$g_W$ scheme map, and a renormalised spectral measure
that can be followed to large $u$.  On a Balaban route this package must also
close Eq.~\eqref{eq:g19-balaban-gate} for the literal path convention and
preserve reflection positivity through every later gauge-fixing and
large/small-field operation.
